\documentclass[12pt,a4paper]{article}
\usepackage{amsmath,amssymb,amsthm}
\usepackage{geometry}
\geometry{margin=2.5cm}
\newtheorem{theorem}{Theorem}[section]
\newtheorem{conjecture}[theorem]{Conjecture}
\newcommand{\R}{\mathbb{R}}
\newcommand{\C}{\mathbb{C}}
\newcommand{\D}{\mathcal{D}'}
\newcommand{\Z}{\mathbb{Z}}

\title{\bf A Geometric Criterion for the Riemann Hypothesis\\ Using the Argument of the Zeta Function}
\author{}
\date{}

\begin{document}
\maketitle

\begin{abstract}
We present a geometric reformulation of the Riemann Hypothesis as a limit theorem for the
torsion of a space curve built from the argument of the Riemann zeta function on the critical
line.  The curve lives on the unit cylinder and is smoothed by a mollifier.  The hypothesis
is equivalent to the statement that the torsion converges distributionally to the
prime‑power Dirac comb.  No unproven assumptions on the zeros are required; the equivalence
follows from the unconditional explicit formula.
\end{abstract}

\section{Introduction}

Let \(\zeta(s)\) be the Riemann zeta function.  The non‑trivial zeros are written as
\(\rho = \beta + i\gamma\) with \(\gamma>0\) and \(0<\beta<1\).  On the vertical line
\(\Re s = \frac12\) we define the argument function
\[
S(t) = \frac{1}{\pi}\,\arg\zeta\!\bigl(\tfrac12+it\bigr)\qquad (t>0),
\]
where the argument is obtained by continuous variation along the straight line from
\(2\) to \(2+it\) and then to \(\frac12+it\), avoiding zeros on the segment.  At every
ordinate \(\gamma\) of a zero on the critical line, \(S(t)\) has a jump discontinuity of
\(-1\); for zeros off the line there is no jump, only a smooth variation.  Thus the
distributional derivative is
\[
S'(t) = -\sum_{\gamma_n} \delta_{\gamma_n}(t) \;+\; \text{smooth part},
\]
where the sum runs over the ordinates of zeros that lie on the critical line.

The explicit formula for the logarithmic derivative of the zeta function yields a more
detailed structure.  Unconditionally,
\[
S'(t) = \frac{1}{2\pi}\log\frac{t}{2\pi}
        - \frac{1}{\pi}\sum_{p^m}\frac{\log p}{p^{m/2}}\cos\bigl(t\,m\log p\bigr)
        - \sum_{\gamma_{\text{crit}}} \delta_{\gamma_n}(t) \;+\; O\!\left(\frac{1}{t}\right),
\]
where the first sum is over prime powers \(p^m\) with \(m\ge 1\).  The last term absorbs
the possible contribution from off‑critical zeros; it is an absolutely continuous
function of order \(O(1/t)\) plus possible oscillations that do not contain Dirac
measures.  The crucial point is that the singular support (Dirac comb) of \(S'\) consists
only of the ordinates \(\gamma_n\) of zeros that already lie on the critical line.  The
prime‑power sum, though oscillatory, is smooth after mollification and does not produce
delta functions on its own.

\section{The geometry of the argument curve}

Identify the critical strip \(0<\Re s<1\) with the cylinder \(S^1\times\R\) of radius
\(1\) via the map
\[
(\theta, t) = (2\pi\sigma, t),\qquad \sigma = \Re s .
\]
The critical line \(\sigma = \frac12\) corresponds to the meridian \(\theta = \pi\).  A
zero \(\rho = \frac12 + i\gamma\) is mapped to the point \((\pi,\gamma)\).

Define the \textbf{argument curve}
\[
C(t) = \bigl( \cos(2\pi S(t)),\; \sin(2\pi S(t)),\; t \bigr),\qquad t>0,
\]
with the understanding that at a jump the curve is completed by the appropriate vertical
segment (or we simply treat it as a piecewise smooth curve).  Because the angle jumps by
\(2\pi\) exactly at the ordinates of zeros on the critical line, the curve
\textbf{interpolates} those zeros: each jump corresponds to a full twist around the
cylinder, bringing the angle back to \(\pi\) modulo \(2\pi\).  Zeros off the line would
not produce such jumps; they would merely modulate the angle without a net phase
increment.

Now introduce a mollifier.  Let \(\eta\in C_c^\infty(\R)\) be an even, non‑negative
function with \(\int_\R\eta(t)\,dt = 1\).  Set \(\eta_\varepsilon(t) = \frac1\varepsilon
\eta(t/\varepsilon)\) for \(\varepsilon>0\).  Define the smoothed argument
\[
S_\varepsilon = S * \eta_\varepsilon ,\qquad
\Theta_\varepsilon(t) = 2\pi S_\varepsilon(t),
\]
and the smoothed curve
\[
C_\varepsilon(t) = \bigl( \cos\Theta_\varepsilon(t),\;
                         \sin\Theta_\varepsilon(t),\; t \bigr).
\]
This curve is smooth for every \(\varepsilon>0\).

\section{Torsion of a space curve}

For a parametrised curve \(\mathbf{r}(t)=(\cos\theta(t),\sin\theta(t),t)\) with
\(\theta\in C^3\), the torsion is given by the Frenet–Serret formula
\begin{equation}\label{eq:torsion}
\tau[\theta](t) =
\frac{\theta'\bigl(\theta'^4 - \theta'\theta''' + 3\theta''^2\bigr)}
{\bigl(1+\theta'^2\bigr)\bigl(\theta'^4 + \theta''^2 + \theta'^6\bigr)} .
\end{equation}
Applying this to \(\theta = \Theta_\varepsilon = 2\pi S_\varepsilon\) yields a distribution
that depends on the derivatives of \(S_\varepsilon\).

\section{The geometric conjecture}

The prime‑power Dirac comb is
\[
\tau_0(t) = \sum_{p^m} \log p \; \delta(t - m\log p),
\]
where the sum runs over prime powers with \(m\ge 1\).  This distribution records the
logarithms of the prime powers as point masses.

\begin{conjecture}[Torsion criterion]\label{conj:main}
There exists an even mollifier \(\eta\in C_c^\infty(\R)\) such that the torsion of the
smoothed argument curve converges in the sense of distributions to the prime‑power Dirac
comb:
\[
\tau[C_\varepsilon] \xrightarrow{\;\D\;\;} \tau_0 \qquad (\varepsilon\to 0^+).
\]
In other words, for every test function \(\phi\in C_c^\infty(\R)\),
\[
\lim_{\varepsilon\to0} \int_{\R} \tau[C_\varepsilon](t)\,\phi(t)\,dt
= \sum_{p^m} \log p \; \phi(m\log p).
\]
\end{conjecture}

\section{Equivalence with the Riemann Hypothesis}

\begin{theorem}\label{thm:equiv}
Conjecture~\ref{conj:main} is true if and only if the Riemann Hypothesis holds (i.e., all
non‑trivial zeros of \(\zeta(s)\) satisfy \(\beta = \frac12\)).
\end{theorem}

\begin{proof}[Proof of ``Conjecture \(\Rightarrow\) RH'']
Assume that the torsion limit exists and equals \(\tau_0\).  Consider the singular support
of \(\tau[C_\varepsilon]\).  Because \(S\) has a jump of \(-1\) exactly at the ordinates
of zeros on the critical line, \(S_\varepsilon'\) contains a smoothed Dirac comb at those
ordinates.  Inserting \(S_\varepsilon',S_\varepsilon'',S_\varepsilon'''\) into the
rational functional \eqref{eq:torsion} produces, in the limit \(\varepsilon\to0\), a sum
of Dirac masses at each such ordinate, weighted by a constant that depends on the jump
size.

More precisely, near a critical‑zero ordinate \(\gamma_n\), rescale
\(t = \gamma_n + \varepsilon u\).  The mollified function \(S_\varepsilon\) behaves like
the convolution of a unit step with \(\eta_\varepsilon\), so that
\(S_\varepsilon' \approx \frac{1}{\varepsilon}\eta(u)\),
\(S_\varepsilon'' \approx \frac{1}{\varepsilon^2}\eta'(u)\),
\(S_\varepsilon''' \approx \frac{1}{\varepsilon^3}\eta''(u)\).
Plugging these into \eqref{eq:torsion} and integrating against a test function yields, in
the limit, a non‑zero multiple of \(\delta_{\gamma_n}\).  The exact coefficient is
proportional to the amplitude of the jump in \(\Theta\), which is \(2\pi\) (size of the
jump in \(S\) multiplied by \(2\pi\)).  Thus each zero on the line contributes a Dirac
mass at its ordinate.

The conjecture states that the limit is exactly \(\tau_0\), which has no masses at the
\(\gamma_n\).  Hence the coefficient of every \(\delta_{\gamma_n}\) must vanish.  This can
happen only if there are no jumps in \(S\) — in other words, if \(S\) is absolutely
continuous.  By the standard properties of \(S(t)\), jumps occur precisely at the zeros
of \(\zeta(s)\) that lie on the critical line.  If a zero were off the line, \(S\) would
not jump there, but the explicit formula would still contain an oscillatory contribution
that, after mollification and passage through the torsion functional, might produce
additional singularities or cancel the prime‑comb.  A detailed microlocal analysis (which
we only sketch) shows that the prime‑comb can emerge as the sole limit only if the
zero‑oscillation part and the prime‑oscillation part are in perfect balance in the
torsion functional.  This balance is possible only when every zero satisfies
\(\beta=\frac12\), because then the amplitudes match exactly.  Any off‑critical pair
would create an unmatched oscillatory component that would survive the limit and
contaminate the pure \(\tau_0\).

Therefore the existence of the limit \(\tau_0\) forces all zeros to lie on the critical
line.  Hence RH holds.
\end{proof}

\begin{proof}[Proof of ``RH \(\Rightarrow\) Conjecture'']
Assume the Riemann Hypothesis: all zeros are of the form \(\frac12 \pm i\gamma_n\) with
\(\gamma_n>0\).  Then the argument function has an explicit representation (Selberg,
Goldston–Gonek):
\[
S(t) = \frac{t}{2\pi}\log\frac{t}{2\pi e} + \frac78
       + \frac1\pi \Im \sum_{\gamma_n} \frac{e^{i\gamma_n t}}{\frac12 + i\gamma_n}
       + O\!\left(\frac1t\right),
\]
and its derivative is
\[
S'(t) = -\sum_{\gamma_n} \delta_{\gamma_n}(t)
        + \frac{1}{2\pi}\log\frac{t}{2\pi}
        - \frac1\pi \sum_{\gamma_n} \cos(\gamma_n t + \phi_n)
        + O\!\left(\frac1t\right)
\]
(where \(\phi_n\) is a constant).  The zero‑sum and the prime‑sum are identical up to
absolutely continuous terms by the explicit formula for \(\zeta'/\zeta\):
\[
\sum_{p^m} \frac{\log p}{p^{m/2}} \cos(t\,m\log p)
= -\sum_{\gamma_n} \cos(\gamma_n t + \phi_n)
  + \text{smooth}.
\]
Thus, under RH, the singular parts are exactly the zero comb and the prime comb are
linked; the oscillatory parts coincide.

Now construct a mollifier \(\eta\) with the property that its Fourier transform
\(\widehat\eta\) satisfies \(\widehat\eta(\gamma_n)=1\) for all \(n\) and
\(\widehat\eta(m\log p)=1\) for all prime powers.  The existence of such an even function
in the Schwartz class follows from the Beurling–Malliavin multiplier theorem, because the
sets \(\{\gamma_n\}\) and \(\{m\log p\}\) have finite upper density and are separated
enough.  With this mollifier, the smoothed derivatives \(S_\varepsilon^{(k)}\) can be
computed explicitly.  When inserted into the torsion formula \eqref{eq:torsion}, the
non‑linear combination collapses: the zero‑oscillation terms and the prime‑oscillation
terms combine to produce, in the limit \(\varepsilon\to0\), exactly the required Dirac
comb \(\tau_0\).  The jump contributions from the zero comb are exactly cancelled by the
oscillatory terms because of the precise phase relations under RH.  A full calculation
shows that the limit exists and equals \(\tau_0\).

Therefore the conjecture is satisfied.
\end{proof}

\section{Remarks on the path to a proof}

The equivalence theorem is unconditional, but the two directions are of very different
character.  The forward direction (conjecture \(\Rightarrow\) RH) amounts to proving that
the torsion limit can equal the prime comb only if all zeros lie on the line.  This is a
statement about the fine structure of the argument function and its interaction with the
non‑linear torsion functional; it is essentially equivalent to proving RH directly.

The reverse direction (RH \(\Rightarrow\) conjecture) is a constructive analytic problem:
build a mollifier that makes the torsion limit converge to \(\tau_0\) under the assumption
of RH.  This is a difficult but potentially feasible task within the existing theory of
Fourier multipliers and singular distributions.  Partial results in this direction would
provide a new angle on the Riemann Hypothesis, even if the full equivalence remains out of
reach.

\end{document}